\documentclass[11pt,letterpaper]{article}
\usepackage[margin=1in]{geometry}
\usepackage{amsmath,amssymb,amsthm,mathtools,microtype}
\usepackage[hidelinks,bookmarksnumbered]{hyperref}
\hypersetup{pdftitle={Integer orbits and reciprocal Mersenne sums},pdfsubject={Irrationality, squarefree semiprimes, and Erdos problem 257}}
\newtheorem{theorem}{Theorem}[section]
\newtheorem{proposition}[theorem]{Proposition}
\newtheorem{corollary}[theorem]{Corollary}
\newtheorem{lemma}[theorem]{Lemma}
\newcommand{\Z}{\mathbb Z}
\newcommand{\Q}{\mathbb Q}
\newcommand{\N}{\mathbb N}
\newcommand{\E}{\mathbb E}
\newcommand{\ind}{\mathbf 1}
\newcommand{\dist}{\operatorname{dist}}
\newcommand{\Var}{\operatorname{Var}}
\newcommand{\Prob}{\mathbb P}
\newcommand{\ee}{\mathrm e}
\newcommand{\eps}{\varepsilon}
\newcommand{\TT}{\mathbb T}
\setcounter{tocdepth}{1}
\setlength{\emergencystretch}{1.8em}
\title{\textbf{Integer orbits and reciprocal Mersenne sums}\\[4pt]
\large Irrationality criteria, prime-factor layers, and the scope of Erd\H{o}s problem 257}
\author{}
\date{15 November 2026}
\begin{document}
\maketitle
\begin{abstract}
A nearest-integer map encodes polynomial relations of a real number as bounded integer sequences. For a real algebraic integer we determine the support and frequency of every defect in its minimal-polynomial recurrence. For reciprocal Mersenne sums we study which infinite sets of exponents force irrationality. One argument gives joint linear independence for prime sums in all ten reduced residue classes modulo eleven at all positive integer dilations. A separate Chinese remainder argument preserves irrationality under arbitrary power-sparse changes to periodic supports. Every infinite support admits a partition into two infinite sets with transcendental sums. We obtain pairwise irrational differences among all finite-prime exclusions of the semiprime sum, and, under Pratt's quantitative prime tuples conjecture, joint independence for the first 69 squarefree layers. The individual full semiprime sum and the universal infinite-support conjecture remain unresolved here.
\end{abstract}
\tableofcontents
\newpage
\section{What the integer orbit determines}
For real $\alpha$ let $r(n)=\lfloor\alpha n+\tfrac12\rfloor$ on $\Z$.
Write $r^0(n)=n$, with superscripts denoting composition. For any
$P(X)=\sum_{k=0}^dc_kX^k\in\Z[X]$ put
$\Delta_P(n)=\sum_{k=0}^dc_kr^k(n)$. Set $e_j(n)=r^j(n)-\alpha r^{j-1}(n)$.

\begin{theorem}[Bounded defect]\label{orbitbound}
For all integers $n$,
\[
 |\Delta_P(n)-nP(\alpha)|\le B_P,\qquad
 B_P=\tfrac12\sum_{k=1}^d|c_k|\sum_{i=0}^{k-1}|\alpha|^i.
\]
Consequently $\Delta_P$ is bounded exactly when $P(\alpha)=0$.
If $P(\alpha)\ne0$, then $|\Delta_P(n)|\to\infty$ as $n\to+\infty$.
\end{theorem}
\begin{proof}
Substitution in $r^k(n)=\alpha r^{k-1}(n)+e_k(n)$ gives
$r^k(n)=\alpha^kn+\sum_{j=1}^k\alpha^{k-j}e_j(n)$.
Use $|e_j(n)|\le1/2$ and sum with coefficients $c_k$.
\end{proof}
In particular $\alpha=p/q$ if and only if $qr(n)-pn$ is bounded for
$q\ge1$. The smallest degree of a nonzero $P$ with bounded defect equals
the degree of $\alpha$ when $\alpha$ is algebraic. When there is no such
$P$, $\alpha$ is transcendental. A single $|\Delta_P(n)|>B_P$ rules out
the particular polynomial. Certifying irrationality this way requires
a proof covering all integer linear polynomials simultaneously.

\subsection*{An exact law for the recurrence defect}
Suppose $\alpha$ is a real algebraic integer of degree $d\ge2$, with
minimal polynomial $P(X)=X^d-\sum_{k<d}c_kX^k$. Write
\[
 P(X)/(X-\alpha)=\sum_{m=0}^{d-1}q_mX^m,\qquad q_{d-1}=1,\qquad
 \sigma=\sum_{m=0}^{d-2}|q_m|.
\]
\begin{theorem}[Support and frequencies]\label{defectlaw}
For $n$ running through the positive integers, the frequency of the defect
$\Delta_P(n)=j$, $j\in\Z$, is
\[
 \operatorname{vol}_{d-1}\!\left\{x\in(-\tfrac12,\tfrac12)^{d-1}:
  \left|\sum_{m=0}^{d-2}q_m x_{m+1}-j\right|<\tfrac12\right\}.
\]
The integers attained for $n\in\Z$ are exactly those with
$|j|<(\sigma+1)/2$. Every attained value has positive frequency for
$n>0$. In particular $r^d=\sum_{k<d}c_kr^k$ on $\Z$ exactly when
$\sigma\le1$.
\end{theorem}
\begin{proof}
The expansion in Theorem~\ref{orbitbound}, or coefficient comparison after multiplying
by $X-\alpha$, gives $\Delta_P(n)=e_d(n)+\sum_{m=0}^{d-2}q_m e_{m+1}(n)$.
Every error has absolute value strictly less than $1/2$, since an integer multiple of
the irrational $\alpha$ cannot be a half-integer. The integer $\Delta_P(n)$ is
therefore the unique nearest integer to the latter sum without $e_d$.

The tuple $(n\alpha,\ldots,n\alpha^{d-1})$ modulo one is equidistributed
on the $(d-1)$-torus by Weyl's criterion: each nontrivial integer
linear combination of $\alpha,\ldots,\alpha^{d-1}$ is irrational, by
minimality. To see what this gives for the errors, put $s(t)=\lfloor t+1/2\rfloor-t$.
At coordinate $j$ the new error is
\[
 e_j=s\!\left(n\alpha^j+\sum_{i<j}\alpha^{j-i}e_i\right).
\]
Given the earlier coordinates, the shift in the argument is fixed, so
the new coordinate sends uniform measure to uniform measure on the
centered unit interval. Iteration gives product measure. These maps are
piecewise affine with discontinuity sets of measure zero, which justifies
using Weyl equidistribution for the slab in the statement. Its boundary
has measure zero because $q_0\ne0$ (the constant coefficient of $P$ is
nonzero). The linear form takes every value in $(-\sigma/2,\sigma/2)$
on the open cube. Its intersection with the open rounding interval for $j$
has a preimage of positive volume exactly under the stated strict inequality.
\end{proof}
For degree two, $\sigma=|\alpha'|$ with $\alpha'$ the conjugate.
Equality to one is impossible for an irrational real quadratic conjugate.
Thus for $\alpha>1$ the exact recurrence holds precisely for quadratic Pisot numbers.
For $\alpha=(1+\sqrt5)/2$ and $1+\sqrt2$, the respective recurrences are
$r^2=r+\mathrm{id}$ and $r^2=2r+\mathrm{id}$. For the real root $\alpha>3$ of
$X^3-3X^2-1$, the two coefficients below the leading term of the quotient
are $1/\alpha$ and $1/\alpha^2$, so $r^3=3r^2+\mathrm{id}$.
For $\alpha=\sqrt2$, $r^2(n)-2n$ equals zero with frequency $1/\sqrt2$,
and each of $-1,1$ with frequency $(1-1/\sqrt2)/2$.

Almost additive models of the reals and identities for iterated floor maps
are treated by Arthan and Fraenkel \cite{arthan,fraenkel}. Kimberling
studies higher-order recurrence partitions, and Shallit studies the
Stolarsky array \cite{kimberling,shallit}. We use a particular nearest-integer tie
convention and give a direct proof of the distribution above. No
priority claim depends on a formula-by-formula identification with these works.

\subsection*{Two finite certificates}
For $\pi$, $n=4{,}000{,}000$ gives $r(n)=12{,}566{,}371$ and
$113r(n)-355n=-77$. Since $77>113/2$, this excludes $355/113$.
For example the enclosing rationals $3.14159265<\pi<3.14159266$ follow
from Machin's identity and the alternating series for both arctangents
using eight terms each.

For Euler's constant $\gamma$ a single input $n=10158$ rules out all
$25^3-1=15{,}624$ nonzero triples $(a,b,c)\in[-12,12]^3\cap\Z^3$:
none satisfies $a+b\gamma+c\gamma^2=0$. Indeed the rational enclosure
$577215/10^6<\gamma<577216/10^6$ forces the successive iterates
$5863,3384$. One may verify the enclosure with $N=1024$, $H_N=\sum_{k=1}^N1/k$,
and the Euler--Maclaurin bounds
\[
 0<\gamma-(H_N-\log N-1/(2N))<1/(12N^2).
\]
For $\log N=10\log2$, take the first ten terms of
$\log2=2\sum_{k\ge0}3^{-2k-1}/(2k+1)$; the remainder is less than
$2\cdot3^{-21}/[21(1-1/9)]$. The complete integer computation is
\begin{verbatim}
min(2000000*abs(10158*a + 5863*b + 3384*c)
    - 1000000*abs(b) - 1577216*abs(c)
    for a in range(-12,13) for b in range(-12,13)
    for c in range(-12,13) if (a,b,c) != (0,0))
\end{verbatim}
Its value $7{,}268{,}352$ is positive. The bound tested uses the
upper endpoint for $\gamma$. The originally reported orbit trials
(12 named numbers, 60 random polynomials, and three near-boundary cases
set aside) come without a full list of inputs; the proofs above are
independent of those trials. No statement about the irrationality of
$\gamma$ follows from this finite certificate. For comparison, the
linear independence of $1,\zeta(2),L(2,\chi_{-3})$ uses arithmetic
holonomy in the work of Calegari, Dimitrov and Tang \cite{cdt}.

\section{Reciprocal Mersenne sums and a first unconditional class}\label{sec:lambert}
For $C\subseteq\N=\{1,2,\ldots\}$, write
\[
 w_m=(2^m-1)^{-1},\qquad S_C=\sum_{c\in C}w_c,\qquad
 T_F=\sum_{\substack{p<r\ \mathrm{primes}\\p,r\notin F}}w_{pr}
\]
for any finite set $F$ of primes. In particular $T_\varnothing$ uses two
\emph{distinct} prime factors. We write $\mathcal S$ for this set of exponents and
$T=T_\varnothing$. Erd\H{o}s problem 257 asks if every infinite $C$
gives irrational $S_C$ \cite{eg}. Absolute convergence holds for every support.
The connection to the preceding orbit test is exact: once an independent
argument proves $S_C$ irrational, each integer linear defect of its
nearest-integer map is unbounded. The existence of the test alone does not
establish irrationality for a specified support.

Write $f_C(m)=\sum_{d\mid m}\ind_{d\in C}$, $E=S_{\N}$ and $\tau=f_{\N}$.
Expansion of the geometric series gives
\begin{equation}\label{lambert}
 S_C=\sum_{m\ge1}2^{-m}f_C(m),\qquad
 T_C(n):=\sum_{h\ge1}2^{-h}f_C(n+h)\equiv2^n S_C\pmod{\Z}.
\end{equation}
Here $T_C(n)$ has an integer argument, whereas $T_F$ without an
argument denotes the prime-exclusion sum above. The congruence follows
by subtracting the terms with $m\le n$. In particular
$f_{\mathcal S}(m)=\binom{\omega(m)}2$, where $\omega(m)$ counts distinct prime factors.


\begin{lemma}\label{finite}
If $B$ is infinite and $\sum_{b\in B}1/b<\infty$, then $S_B$ is irrational.
\end{lemma}
\begin{proof}
The periodic weight $w_b(n)=2^{n\bmod b}/(2^b-1)$ satisfies $T_B(n)=\sum_{b\in B}w_b(n)$ by summing $2^{-h}$ over $b\mid n+h$. For $n\ge1$, $T_B(n)-S_B>0$ since each difference $w_b(n)-w_b(0)$ is nonnegative and those with $b>n$ are strictly positive. If $qS_B\in\Z$, \eqref{lambert} would thus give $T_B(n)-S_B\ge1/q$ for every $n\ge1$.

Let $D_0$ be fixed and take means over $n=D_0j$, $1\le j\le J$. With $g=(D_0,b)$, the mean of $w_b(D_0j)$ over one period is $g/[b(2^g-1)]\le 1/b$. For all $J$, the finite mean is bounded by $(D_0+2)/b$: for each $h\ge1$ there are at most $(D_0J+h)/b$ multiples of $b$ among $D_0j+h$, so one can sum with weights $2^{-h}$. Dominated convergence permits termwise passage to $J\to\infty$. Now let $D_0=\operatorname{lcm}(1,\ldots,m)$ and let $m\to\infty$. For each $b$ the limiting mean eventually equals $w_b(0)$, with common majorant $1/b$. A second dominated convergence shows that the limiting means of $T_B(n)-S_B$ tend to zero, contradicting their lower bound. This averaging proof is also recorded in \cite{cook}.
\end{proof}

\section{Transcendental subseries inside every infinite support}\label{sec:lacunary}
For this section no estimate for primes or divisor correlations is needed.
Our starting observation is that a simultaneous rational approximation
which is too accurate for the degree of a polynomial forces an
\emph{exact} polynomial relation at each sufficiently late truncation.

\begin{theorem}[Sparse independent coordinates]\label{thm:algfamily}
Let $n_1<n_2<\cdots$ be positive integers, and let
$D_k=\operatorname{lcm}_{i\le k}(2^{n_i}-1)$. Assume
\[
 \frac{n_{k+1}}{\log_2(2D_k)}\longrightarrow\infty.
\]
Let $\mathcal I$ be any family of infinite subsets of $\N$ such that
each two distinct members have finite intersection. Then the family
\[
 X_I=\sum_{k\in I}\frac1{2^{n_k}-1}\quad (I\in\mathcal I)
\]
is algebraically independent over $\Q$. In particular every member
is transcendental.
\end{theorem}
\begin{proof}
Consider a finite selection of distinct sets $I_1,\ldots,I_r$, put
$x=(X_{I_1},\ldots,X_{I_r})$, and truncate all coordinates to $k$ to
obtain $x^{(k)}$. Each coordinate has denominator dividing $D_k$.
Moreover $\|x-x^{(k)}\|_\infty\le2^{2-n_{k+1}}$ by the geometric tail.
If a nonzero polynomial vanishes at $x$, take one of least total degree,
with integer coefficients after clearing denominators; denote it by $F$
and its degree by $m$. On the compact box containing the truncations
its first derivatives are bounded. Hence $|F(x^{(k)})|\ll_F2^{-n_{k+1}}$.
This is smaller than $D_k^{-m}$ for all sufficiently large $k$.
Since $D_k^m F(x^{(k)})$ is an integer, $F(x^{(k)})=0$ eventually.

For every $i$ there are infinitely many $k+1$ belonging to $I_i$
and to none of the other $r-1$ sets. At those steps $x^{(k+1)}$
changes only in coordinate $i$, by the strictly positive number
$(2^{n_{k+1}}-1)^{-1}$. Both values of $F$ are zero. The one-coordinate
difference quotient therefore has limit $\partial_iF(x)=0$.
A nonconstant polynomial in characteristic zero has a nonzero formal
partial derivative. This gives a vanishing polynomial of smaller degree,
including a contradiction if that derivative is constant.
\end{proof}

\paragraph{Ten explicit numbers.}
For $j=0,\ldots,9$ put
\[
 X_j=\sum_{\substack{k\ge1\\v_2(k)=j}}\frac1{2^{(k+3)!}-1}.
\]
The factorials form a divisibility chain, so $D_k=2^{(k+3)!}-1$
exactly. The ratio in Theorem~\ref{thm:algfamily} grows like $k+4$.
The ten sets of indices are infinite and disjoint. Thus these ten
specified numbers are jointly algebraically independent. Every
nonzero integer polynomial in any one $X_j$ yields an unbounded orbit
defect by Theorem~\ref{orbitbound}. Polynomial approximation supplies
the global argument here. This construction does not assign a new
status to ten previously identified open constants.

\begin{corollary}[A conclusion for every infinite support]\label{cor:arbitrarysplit}
For every infinite $A\subseteq\N$ there is a continuum-sized family of
infinite subsets $B_\xi\subset A$ for which all $S_{B_\xi}$ are jointly
algebraically independent. Every two distinct complementary sums
$S_{A\setminus B_\xi}$ have transcendental difference. At most one of
these complementary sums can be algebraic. In particular every infinite
support can be partitioned into two infinite supports whose sums are both
transcendental.
\end{corollary}
\begin{proof}
Choose a subsequence $n_k$ of $A$ successively with
$n_{k+1}>(k+1)(1+\sum_{i\le k} n_i)$, skipping at least one element of
$A$ between choices. The inequality
$\log_2D_k\le\sum_{i\le k}n_i$ implies the theorem's hypothesis. For each
infinite binary word $\xi$ let $I_\xi$ consist of the positive integers
whose binary digits are a leading $1$ followed by a nonempty finite prefix
of $\xi$. Two different words have finitely many common prefixes. Take
$B_\xi=\{n_k:k\in I_\xi\}$. Alleged algebraicity of
$S_{B_\xi}-S_{B_\eta}$ for $\xi\ne\eta$ would, by composition with
the minimal polynomial of the difference, contradict algebraic
independence. The complementary sums have the opposite differences,
so two cannot both be algebraic. The skipped elements keep every
complement infinite.
\end{proof}

Apply Corollary~\ref{cor:arbitrarysplit} with $A=\mathcal S$, the actual
squarefree semiprimes. One can choose $n_k$ deterministically as the
least allowed semiprime beyond the indicated bound and beyond the next
unused exponents. The supports $\mathcal S\setminus B_\xi$ consist
entirely of semiprimes and retain every semiprime outside a single
rapidly growing sequence. Among their continuum many sums at most
one is algebraic. The full sum $T$ is common to every identity
$S_{\mathcal S\setminus B_\xi}=T-S_{B_\xi}$ and cancels from every
difference. If $T$ is algebraic, all of these complementary sums are
jointly algebraically independent: algebraic independence persists
over $\Q(T)$, and coordinatewise translation and change of sign are
invertible there. This implication imposes no contradiction on the
hypothesis that $T$ is algebraic.

\begin{theorem}[An explicit partition for every infinite support]\label{thm:all-but-full}
For any infinite $A\subseteq\N$ and any fixed $s\ge2$ there is an
explicit partition into $s$ infinite sets $A_0,\ldots,A_{s-1}$ such
that for every nonempty proper $J\subset\{0,\ldots,s-1\}$ the number
\[
 S_{\bigcup_{i\in J} A_i}=\sum_{i\in J}S_{A_i}
\]
is a Liouville number. In particular it is transcendental.
\end{theorem}
\begin{proof}
Set $M_0=8$ and, for $j\ge0$, take
\[
 M_{j+1}=\max\big(2^{M_j},\,\min\{a\in A:a>M_j\}\big).
\]
Assign $A\cap(M_j,M_{j+1}]$ to $A_{j\bmod s}$, and assign
$A\cap[1,M_0]$ to $A_0$. Every block contains an element, so each
part is infinite. For fixed nonempty proper $J$ put $x=\sum_{i\in J}S_{A_i}$.
There are infinitely many $j$ with $j\bmod s\notin J$. Truncate $x$ at
$M_j$ to a reduced rational $a_j/b_j$ at these indices. The entire
next block is absent, and later selected blocks contain positive terms.
The denominator divides $\prod_{n=1}^{M_j}(2^n-1)$, so
\[
 0<x-a_j/b_j\le2^{1-M_{j+1}},\qquad
 \log_2 b_j\le M_j(M_j+1)/2.
\]
Along the indicated indices the truncations converge to $x$ and
take infinitely many distinct values, because further selected blocks
occur. Their reduced denominators tend to infinity. For every fixed
integer $h$ the displayed error is eventually less than $b_j^{-h}$.
This also rules out rational $x=a/b$, which would give the lower
bound $1/(bb_j)$. Thus $x$ is Liouville. For completeness, for a
real algebraic irrational of degree $m$ with integer minimal polynomial
$P$, $|P(a_j/b_j)|\ge b_j^{-m}$, while boundedness of $P'$ near $x$
gives $|P(a_j/b_j)|\ll_P|x-a_j/b_j|$. Taking $h>m$ is impossible.
\end{proof}

For $A=\mathcal S$ and $s=10$ one may simply take $M_0=8$,
$M_{j+1}=2^{M_j}$. Bertrand's postulate supplies an odd prime
$p$ with $M_j<2p<2M_j$ for these even $M_j$, so every block
contains a squarefree semiprime. The ten fixed sums consist entirely of
terms of the full semiprime sum. All $2^{10}-2=1022$ nonempty proper
unions have Liouville sums. They are distinct by strict tail
domination (proved below). Their ten individual sums add up to $T$.
For the full index set there is no omitted block, so the displayed
approximation at $M_j$ no longer follows. In particular this gives an
explicit partition of the \emph{original} target into ten provably
transcendental summands without deciding the arithmetic nature of their total.


\section{Prime classes at different dilations}

For a set $\mathcal P$ of primes put
$\omega_{\mathcal P}(m)=\sum_{p\in\mathcal P,\ p\mid m}1$.
The next result permits different periodic weights on different prime
sets, together with signed contributions from the composite exponents.

\begin{theorem}\label{mixcube:main}
Let $\mathcal P_1,\ldots,\mathcal P_s$ be disjoint arbitrary sets of
primes and let $w_1,\ldots,w_s$ be nonzero periodic integer-valued
functions with common period $D$. Assume
$\sum_{p\in\bigcup_i\mathcal P_i}1/p=\infty$. Let $\mathcal C$ be a
set of composites such that each prime divides at most $B$ members,
where $B$ is fixed, and let $\lambda_c$ be bounded integers. Then
\[
 \sum_{m\ge1}2^{-m}\sum_{i=1}^s w_i(m)\omega_{\mathcal P_i}(m)
       +\sum_{c\in\mathcal C}\frac{\lambda_c}{2^c-1}
 \quad\hbox{is irrational.}
\]
\end{theorem}

\begin{proof}
We adapt the affine prime cube of Tao and Ter\"av\"ainen \cite[Section~5]{tt} to separate periodic weights and composite divisors. Write
$\mathcal P=\bigcup_i\mathcal P_i$, and set $w_p=w_i$ on $\mathcal P_i$
and $w_p=0$ on primes outside $\mathcal P$. Put
\[
 f(m)=\sum_{p\mid m}w_p(m)+\sum_{c\in\mathcal C,\ c\mid m}\lambda_c .
\]
The displayed number is $\sum_m2^{-m}f(m)$. If $q$ times this number
is an integer, then $q\sum_{j\ge1}2^{-j}f(k+j)$ is an integer for
all $k\ge0$. Write $e(x)=\exp(2\pi i x)$ and
$t=\sum_{p\in\mathcal P,\,p\le3X}1/p$, $z=\log(2+t)$,
$u=\log\log X$. Choose $X\to\infty$
so that $\sum_{p\in\mathcal P,\,R<p\le3X}1/p\le z^{3/2}$ for
$R=X^{1/u^2}$. Indeed, the nondecreasing function
$m(u)=\sum_{p\in\mathcal P,\,p\le3\exp(\exp u)}1/p$ is at most $u+O(1)$.
If its increment on $[u-3\log u,u]$ exceeded $[\log(2+m(u))]^{3/2}$
at every sufficiently large $u$, summing the increments across $[u/2,u]$
would first give $m(u)\gg u/\log u$ and then $m(u)\gg u\sqrt{\log u}$.
The resulting contradiction gives the required subsequence since
$R\ge3\exp(\exp(u-3\log u))$ eventually.

Take $K=\lfloor2\log z/\log 2\rfloor$, $H=\lceil z^2\rceil$,
$P=\lceil\exp(z^3)\rceil$, and $H_+=\lceil u^2\rceil$. Take $\epsilon\in\{0,1\}^K$ and $1\le j\le K$ for the steps. Use an
affine cube of distinct auxiliary primes $p_\epsilon\in[P/2,P]$ all
congruent to $1\pmod D$, with
$p_\epsilon=p_0+\sum_j\epsilon_jv_j$, $t_\epsilon=\sum_j j\epsilon_jv_j$,
$r_{\epsilon,j}=p_\epsilon j-t_\epsilon$, and distinct near shifts
$a_{\epsilon,h}=r_{\epsilon,K+h}$, $1\le h\le H$.
For completeness, embed the eligible primes in the field of prime
order $T\in[5P,10P]$. Their density is $\gg_D1/\log P$ by the prime
number theorem in fixed arithmetic progressions. Successive
Cauchy--Schwarz inequalities (monotonicity of the Gowers norms) give
probability at least $(c_D/\log P)^{2^K}$ that a uniformly chosen affine
cube is contained in this set. A collision between two vertices or
two distinct near shifts imposes a nonzero linear equation. Two
different $h$ have different coefficients of $p_0$, whereas for equal
$h$ a differing bit $j$ has coefficient difference $\pm(K+h-j)$.
All these coefficients are smaller than $T$ in absolute value. Thus
the total collision probability is $O(2^{2K}H^2/T)$, smaller than the
displayed cube probability. Lift vertices to their integer prime
representatives. Each edge difference is in $[-P/2,P/2]$, so parallel
edges have the same integer difference by uniqueness of this lift.
The cube therefore lifts with $|v_j|\le P/2$ and $D\mid v_j$. Write $W$ for
the product of the absolute differences of distinct near shifts. Set
$Q_0=\operatorname{lcm}(D,\operatorname{rad}(W\prod_\epsilon p_\epsilon))$.
There are $O(2^{2K}H^2)$ factors in $W$, each
$O((H+K)KP)$, and thus $\log Q_0\ll_D (1+z)^{O(1)}$.
Fix a class modulo $Q_0$ satisfying $n\equiv b-K-1\pmod D$ for any
fixed $b$, $n\equiv t_\epsilon\pmod{p_\epsilon}$, and $n\equiv0\pmod p$
at the remaining primes outside $D$. Then $D\mid v_j$ and all
$p_\epsilon$ occur to the first power in $Q_0$.

Write $U=(2+P Q_0+z)^3$ and freeze the set
$\mathcal F=\{c\in\mathcal C:c\le z\}\cup
\{c\in\mathcal C:c\le U,\ (c,Q_0)>1\}$ on lifts of the coarse class
modulo $Q_1=Q_0\prod_{c\in\mathcal F}c$. We have
$\log Q_1\ll_{B,D}(1+z)^{O(1)}$, using
$\#\{c\in\mathcal C:c\le y\}\le B\sqrt y$ and the fact that at most
$B\omega(Q_0)$ composites meet $Q_0$. The first bound follows by
charging a composite to its least prime factor. For $N\in[\sqrt X,X]$
let coarse and lift averages run over $N<n\le2N$. For any shifted
argument $n+r\in(0,3N]$, the coarse expected number of divisors in
$\mathcal C\setminus\mathcal F$ is
\[
 O_B\big(z^{-1/2}+\omega(Q_0)Q_0/U+Q_0/\sqrt N\big)=o(1).
\]
Indeed for each $c$ the probability is at most $(c,Q_0)/c+O(Q_0/N)$.
Dyadic summation bounds the reciprocal tail of all composites by
$O(Bz^{-1/2})$, and any remaining composite meeting $Q_0$ exceeds $U$.

For an auxiliary prime $p=p_\epsilon$ and $k=(n-t_\epsilon)/p+j$,
we have $pk=n+r_{\epsilon,j}$ and $pk\equiv k\equiv n+j\pmod D$.
The exact dilation difference is
\[
 f(pk)-f(k)=w_p(n+j)(1-\mathbf1_{p\mid k})
   +\sum_{c\in\mathcal C,\,p\mid c}\lambda_c
     (\mathbf1_{c/p\mid k}-\mathbf1_{c\mid k}).
\]
The first term $w_p(n+j)$ is deterministic, even when different cube
vertices belong to different prime sets. Its infinite weighted sum
can be absorbed in a constant phase. The exceptional event $p\mid k$
has probability at most $1/p+O(Q_0/N)$ since $p^2\mid n+r_{\epsilon,j}$.
For $c\in\mathcal F$ the dilation corrections are constant on each lift,
since $pc\mid Q_1$. At most $B$ remaining composites occur for this
vertex. The correction for each is supported on $c\mid n+r_{\epsilon,j}$
and has coarse probability $O(Q_0/U+Q_0/N)$. Thus, after summing the
dilation errors over all vertices and tail indices with weights
$2^{-j}$, their coarse mean absolute size is $O(2^K(P^{-1}+Q_0/U+Q_0/N))$.

Use the alternating signs $\sigma_\epsilon=(-1)^{\sum_j\epsilon_j}$.
For $j\le K$, the argument $n+r_{\epsilon,j}$ is independent of
$\epsilon_j$, so the values of the entire $f$ cancel. For the remaining
$j=K+h$ the total absolute weight is at most one. The coarse composite
estimate discards the unfrozen composites for $h\le H_+$, and beyond
$H_+$ the bound $|f(m)|\ll_{B,w,\lambda}\log(2m)$ suffices. For
the prime part the tail $H<h\le H_+$ has coarse weighted mean absolute
size $O(2^{-H}(1+t+\omega(Q_0)+X^{1/100}Q_0/N))$, by counting primes
up to $X^{1/100}$ one at a time and using the bounded number of larger
divisors of each shifted argument. These errors all tend to zero.
At a near shift the weight for color $i$ is
$c_{i,h}=w_i(b+h-1)$, independent of the vertex. Define, for
$p\in\mathcal P_i$,
\[
 T_p(n)=q\sum_{h=1}^H\sum_\epsilon 2^{-K-h}\sigma_\epsilon c_{i,h}
                         \mathbf1_{p\mid n+a_{\epsilon,h}}.
\]
Primes dividing $Q_1$ contribute a constant on each lift. The integral
tail identities now give a coarse average
$\mathbb E[\eta(n)e(\sum_{p\in\mathcal P,\ p\nmid Q_1}T_p(n))]=1+o(1)$,
where $\eta$ is constant of modulus one on each lift. All errors so far
were estimated before conditioning on a lift.

For independent uniform residues $V_p\pmod p$ and $p\nmid Q_1$ in color
$i$, the near shifts are distinct modulo $p$ and hence
\[
 \mathbb E\,T_p(V_p)=0,\qquad |T_p|\ll 2^{-K},\qquad
 \mathbb E\,T_p(V_p)^2=\frac{q^2}{2^Kp}\sum_{h\le H}4^{-h}c_{i,h}^2.
\]
The last sum is bounded below by $4^{-D}$ for $H\ge D$, since the
first $D$ arguments exhaust all residues and $w_i$ is integer valued
and nonzero. It is bounded above independently of the color. The sum
of variances for $p\le R$, $p\nmid Q_1$ is consequently
$V\asymp t/2^K\to\infty$. Indeed the annulus costs $O(z^{3/2})$ in
reciprocal mass and primes dividing $Q_1$ cost $O(1+\log z)$. Also
$V=o(u)$. For a local variable $T$ and an independent copy $T'$, one has
$1-|\mathbb E e(T)|^2=\mathbb E(1-\cos(2\pi(T-T')))\gg\operatorname{Var} T$
since $|T|=o(1)$. Multiplication gives $O(\exp(-cV))$ for the
absolute independent phase mean. On every lift moments through
$M=2\lceil u\rceil$ agree with those of the independent model up to
$O(R^{2M}Q_1/N)$, by expanding into at most $R^M$ terms of supremum at most one
and period at most $R^M$ coprime to $Q_1$. The complete-period average
of each term equals its independent average by the CRT, even with
repeated prime indices, and the incomplete-period error is $O(R^MQ_1/N)$.
Here $R^{2M}Q_1/N=X^{-1/2+o(1)}$ as an upper bound. The independent centered variables give exponential moments
at the two fixed parameters $\pm2\pi\mathrm e$ bounded by $\exp(O(V))$.
Thus the Taylor remainder for the characteristic function is
$O(\exp(O(V)-M))+o(1)$ by the even moment of order $M$ and the moment
comparison. The medium-prime phase average tends to zero uniformly on
all lifts. For $R<p\le Y=X^{1/100}$ the sum has lift mean square
$O(Y^4Q_1/N+2^{-K}z^{3/2})=o(1)$, by complete-period orthogonality
for distinct prime variables.

Finally split the prime bins above $Y$ by color. With
$L=\lceil(\log X)^{1/\log u}\rceil$, this results in $O(sL)$ bins
of harmonic mass $O(1/L)$. For each color use
$F_i(m)=\sum_{p\in\mathcal P_i,\,Y<p\le3X}\mathbf1_{p\mid m}$.
On each lift its approximation by the sum of the bin-hit indicators
has failure probability $O(Q_1/L)$. There are at most
$3N/(pp')$ possible arguments with two distinct prime divisors $p,p'$
in one bin. Divide by the lift size $\asymp N/Q_1$ and sum using that
the sum of squared bin masses is $O(1/L)$. Each bin avoidance function is completely multiplicative.
The equidistribution test for Theorem~3.1(i) of Tao and Ter\"av\"ainen
\cite{tt} still applies to an arbitrary subset of a bin: the inclusion--
exclusion error on an interval $(T,2T]$, $X^{0.4}\le T\le X$, is
$O(T/\log X)$ for moduli up to $(\log X)^{1/4}$. This counts squarefree
products of primes above $Y$ by removing the largest prime and using
$\prod_{Y<p\le2X}(1+1/p)=O(1)$. Larger moduli satisfy the required
bound trivially, and the technical small-prime condition is automatic.
Thus a common scale may be selected for all $O((sL)^2)$ bin
pairs and bin--constant pairs. The exceptional sets in $[\sqrt X,X]$ have normalized logarithmic measure $O((\log X)^{-c})$ each, hence their union has measure $o(1)$. The resulting estimates are uniform on lifts and near
shifts, since $Q_1$, the shift sizes, and $L$ are $(\log X)^{o(1)}$
and $Q_1/L=o(1)$. At the selected $N$, distinct shifts have centered
$F_i$ covariance $O(Q_1/L+L^2(\log X)^{-c})$ for some $c>0$. In detail, writing $g_J$ for a bin avoidance function and
$\delta_{J,N}=N^{-1}\sum_{N<m\le2N}g_J(m)$, the theorem applies
both to $(g_J,g_{J'})$ and to $(g_J,1)$. Subtract $\delta_{J',N}$
times the latter estimate from the former and sum over both bins in
the given color. This bounds the covariance centered at the same
$C_{i,N}=\sum_{J\text{ in color }i}(1-\delta_{J,N})$ on every lift.
Each $F_i$ and each sum of bin-hit indicators is bounded by the
absolute maximum number of primes exceeding $X^{1/100}$ dividing
one integer $m\le3X$. Taking the unrestricted interval mean therefore
also bounds $C_{i,N}$ up to $O(L/N)$. The alternating
average over vertices at a fixed $h$ therefore has $L^2$ norm $o(1)$,
including the diagonal contribution $O(2^{-K})$ to its second moment.
Multiplying by $q c_{i,h}2^{-h}$ and summing over $h$ and the finitely
many colors gives an $o(1)$ very-large-prime contribution on each lift.
Consequently the total phase has mean $o(1)$ uniformly on all lifts.
Multiplication by $\eta$ and averaging gives the required contradiction.
\end{proof}

\begin{corollary}\label{mixcube:grid}
Let $\mathcal P_i$, $i$ in an arbitrary index set, be pairwise disjoint
sets of primes, each with divergent reciprocal sum. Then the numbers
\[
 1,\qquad V_{i,d}=\sum_{p\in\mathcal P_i}\frac1{2^{dp}-1}
            \quad (d\ge1)
\]
are linearly independent over $\mathbb Q$. In particular take the ten
nonzero prime residue classes modulo $11$. Writing $S_a=V_{a,1}$ and
$L_d=\sum_p(2^{dp}-1)^{-1}$, the family
$1,S_1,\ldots,S_{10},L_2,L_3,\ldots$ is linearly independent.
Each $L_p$ with $p$ prime can instead be represented by the sum $U_p$
over squarefree semiprimes with least prime factor $p$.
\end{corollary}
\begin{proof}
Every finite combination with integer coefficients $b_{i,d}$ differs
by a rational number from the series in Theorem~\ref{mixcube:main}
with $\mathcal C=\varnothing$ and $w_i(m)=\sum_{d\mid m}b_{i,d}$.
Indeed $V_{i,d}=\sum_{d\mid m}2^{-m}\omega_{\mathcal P_i}(m/d)$, and
on multiples of $d$ the change from $\omega_{\mathcal P_i}(m)$ counts
only $p\in\mathcal P_i$ dividing $d$ with $v_p(m)=v_p(d)$. The
exceptional contribution is periodic in $m$ and its binary series is
rational. If a row has a nonzero coefficient, its least such index
$d$ gives $w_i(d)\ne0$. Discard zero rows and apply the theorem. The
ten residue classes have divergent reciprocal sum by Mertens' theorem
for fixed arithmetic progressions \cite{williams}. In this specialization the exact
relations are
\[
 L_d=\sum_{a=1}^{10} V_{a,d}+\frac1{2^{11d}-1},\qquad
 U_p=L_p-\sum_{\substack{r\le p\\r\ {\rm prime}}}\frac1{2^{pr}-1}.
\]
The second formula makes replacement valid in the presence of $1$.
The first explains the restriction $d>1$ when adjoining $L_d$ to the
$S_a$. More generally, if $1$, all $V_{a,d}$, all $L_d$ and all $U_p$
are listed together, these displayed relations generate the entire
space of finite rational relations: eliminate $U_p$ and $L_d$ first,
then apply the linear independence of $1$ and the grid.
\end{proof}

\subsection{Bounded incidence and periodic weights}

\begin{corollary}\label{mixcube:incidence}
If $A$ is infinite and each prime divides at most $B$ composite
members of $A$, for some fixed $B$, then $S_A$ is irrational.
In particular this holds for every infinite pairwise coprime set,
every infinite subset of the primes, and every infinite subset of
$\{p_j:j\ge1\}\cup\{p_jp_{j+1}:j\ge1\}$, where $p_j$ is the $j$th
prime. Every nonzero periodic integer weight $w$ also gives an
irrational sum $\sum_{m\ge1}2^{-m}w(m)\omega(m)$.
\end{corollary}
\begin{proof}
The prime members of $A$ supply a single color with constant weight one, and
its composite members have $\lambda_c=1$. If the prime reciprocal sum
diverges, apply Theorem~\ref{mixcube:main}. If it converges, charging
composites to their least prime factors gives $O(B\sqrt y)$ composites
up to $y$, so the reciprocal sum of the composites converges as well.
Apply Lemma~\ref{finite} to the full $A$. Including $1$
only changes $S_A$ by an integer. In the three examples one may take
$B=1,0,2$, respectively. For the periodic claim use all primes as one
color and take $\mathcal C$ empty.
\end{proof}

For a finite set $F$ of primes write
\[
 T_F=\sum_{\substack{p<r\ {\rm primes}\\p,r\notin F}}(2^{pr}-1)^{-1}.
\]
Then the following identity specifies all finite corrections:
\[
 T_F=T_{\varnothing}-\sum_{p\in F}L_p
       +\sum_{p\in F}\frac1{2^{p^2}-1}
       +\sum_{\substack{p,r\in F\\p<r}}\frac1{2^{pr}-1}.
\]
Indeed, $L_p$ counts one square and all two-prime products containing
$p$. A product with both primes in $F$ is counted twice in the sum.
For $F\ne G$, subtraction gives $T_F-T_G$ equal to a nonzero
integer combination of distinct $L_p$ plus a rational number.
Corollary~\ref{mixcube:grid} therefore makes this difference irrational.
At most one member of the family $(T_F)_F$ can be rational. Also the
$U_p$ have disjoint supports and $T_{\varnothing}=\sum_pU_p$.
Their finite linear independence does not decide the status of this
infinite sum.

\subsection{Joint centered rounding errors}

Choose any distinct finite set of grid entries $x_1,\ldots,x_t$
from Corollary~\ref{mixcube:grid}, including for example the ten
$S_a=V_{a,1}$ for $1\le a\le10$. For $r_x(n)=\lfloor nx+1/2\rfloor$,
\[
 |v r_{x_j}(n)-an|\longrightarrow\infty
 \quad(v\in\mathbb Z\setminus\{0\},\ a\in\mathbb Z).
\]
Indeed the left side is $|n(vx_j-a)+O(|v|)|$ with $vx_j-a\ne0$.
Moreover the vectors $(r_{x_j}(n)-nx_j)_{j=1}^{t}$ are uniformly
distributed with respect to volume on $[-1/2,1/2]^t$. For each
nonzero $k\in\mathbb Z^t$, the number $k\cdot x$ is irrational, so
$N^{-1}\sum_{n\le N}e(nk\cdot x)\to0$ by the finite geometric
sum, where $e(y)=\exp(2\pi i y)$. Weyl's criterion gives uniform
distribution of $nx$ modulo one. Coordinatewise centered rounding is
a measure-preserving piecewise isometry outside finitely many
coordinate hyperplanes of measure zero. The same conclusions hold for
any finite subfamily of $S_1,\ldots,S_{10},L_2,L_3,\ldots$, with the
permitted replacements by $U_p$.


\section{Sparse deletions and a simultaneous return of two tails}
\label{sec:crtreturn}
Write $E=S_{\mathbb N}$ and $\tau(m)=f_{\mathbb N}(m)$. Denote the
scaled tail for support $C$ by $T_C(n)=\sum_{h\ge1}2^{-h}f_C(n+h)$.

\begin{theorem}\label{crt:main}
Suppose that $B\subseteq\N$ and $M_s=\sum_{b\in B}b^{-s}<\infty$ for some $0<s<1$. Then $S_{\N\setminus B}$ is irrational. If $B$ is infinite then $1,S_B,S_{\N\setminus B}$ are linearly independent over $\Q$.
\end{theorem}
For every such $B$, $\#(B\cap[1,x])\le M_s x^s$. Its complement has density one, divergent reciprocal sum and infinitely many composite multiples of every prime. Thus Theorem~\ref{crt:main} treats cases in which no uniform bound on the number of composite multiples of a prime is possible. We use the classical irrationality of $E$ proved by Erd\H{o}s \cite{erdos} and Lemma~\ref{finite}.

\subsection{Returning the sparse tail while forcing a divisor block}
\begin{lemma}\label{crt:return}
Under the hypothesis of Theorem~\ref{crt:main}, there are $n\to\infty$ for which
\[
 \dist(T_{\N}(n),\Z)\longrightarrow0,\qquad T_B(n)\longrightarrow S_B.
\]
\end{lemma}
\begin{proof}
Let $k\to\infty$, $L=k^3$ and $z=\log k$. A positive integer is called $z$-smooth if every prime dividing it is at most $z$; $1$ is included. Set $D=\prod_{p\le z}p^{e_p}$, where $e_p$ is the first positive integer such that $p^{e_p}>2L$. Then $\log D=O((\log k)^2)$. The condition $n\equiv0\pmod D$ gives
\begin{equation}\label{crt:small}
 v_p(n+h)=v_p(h)\quad (p\le z,\ 1\le h\le L).
\end{equation}

Choose an integer $K>\max(4,3/(1-s))$ and put $P=k^K$, $t=k(k+1)/2$. Let $F$ be a uniformly random $t$-element subset of the primes in $[P,2P]$. Their number $J\gg P/\log P$ by the prime number theorem. Define
\[
 W_F(b)=\prod_{\substack{r\in F\\v_r(b)=1}}r .
\]
For any prescribed set of $j$ primes in this interval, its probability of inclusion is at most $(t/J)^j$, as follows from sampling without replacement. Expand $W_F(b)$ as a product of factors $1+(r-1)\ind_{r\in F}$ over primes $r\in[P,2P]$ with $v_r(b)=1$. If there are $m$ such primes, the expansion gives
\[
 \E_F W_F(b)\le (1+2Pt/J)^m\le P^{(1-s)m}\le b^{1-s}
\]
for large $k$, since $2Pt/J\ll k^2\log k$. By nonnegativity we may sum this expectation over all $b\in B$. We may therefore fix a single $F$ for which
\begin{equation}\label{crt:good}
 \sum_{\substack{b\in B\\b\text{ not }z\text{-smooth}}}\frac{W_F(b)}{b}
 \le \sum_{\substack{b\in B\\b>z}} b^{-s}=o(1).
\end{equation}

Assign $i$ primes from $F$ to shift $i$, for each $1\le i\le k$, using each prime once. For every $r$ assigned to $i$, impose $n+i\equiv r\pmod{r^2}$, together with $n\equiv0\pmod D$. These conditions define one class modulo
\[
 Q=D\prod_{r\in F}r^2,\qquad \log Q\asymp k^2\log k.
\]
Hence $2^i\mid\tau(n+i)$ for $1\le i\le k$: each assigned prime has exact valuation one and doubles the divisor count. If $h\le L$, a prime in $F$ has valuation zero in $n+h$ unless it is assigned to $h$, since $r>2L$. In particular all such valuations are at most one.

Fix an integer $C>\max(3,2/(1-s))$. Write $X=Q^C$ and average uniformly over our class in $X<n\le2X$. There are exactly $X/Q$ points. For any positive $d$ the condition $d\mid n+h$, if soluble, has probability at most
\begin{equation}\label{crt:crt}
 \frac{(d,Q)}{d}+\frac QX.
\end{equation}
This is the count of one residue class modulo $d/(d,Q)$ among $X/Q$ consecutive integers.

For $h\le L$ and $b\in B$ not $z$-smooth, if $b\mid n+h$ is possible then the part of $b$ supported on $p\le z$ divides $h$ by \eqref{crt:small}. No prime in $F$ divides $b$ to exponent two or higher. Thus $(b,Q)\le h W_F(b)$ for each soluble condition. With $b\le 3X$ automatically in this sum, \eqref{crt:crt} and the growth estimate for $B$ yield
\[
 \E_n\sum_{\substack{b\in B,\ b\mid n+h\\b\text{ not }z\text{-smooth}}}1
 \le h\sum_{\substack{b\in B\\b\text{ not }z\text{-smooth}}}\frac{W_F(b)}{b}
   +O_s(M_s Q X^{s-1}).
\]
The endpoint error tends to zero by the choice of $C$. For $z$-smooth $b$, \eqref{crt:small} gives $b\mid n+h$ exactly when $b\mid h$. The contribution of nonsmooth divisors of $h$ to the original tail is bounded by $\sum_{h>z}2^{-h}\tau(h)=o(1)$. Since $\sum h2^{-h}=2$, \eqref{crt:good} now shows that
\begin{equation}\label{crt:bshort}
 \E_n\sum_{h\le L}2^{-h}|f_B(n+h)-f_B(h)|=o(1).
\end{equation}
For $h>L$ we may use $f_B(n+h)\le\tau(n+h)\le2X+h$ and $\tau(h)\le h$. The tail costs $O((X+L)2^{-L})=o(1)$ because $\log X=O_s(k^2\log k)$. Consequently $\E_n|T_B(n)-S_B|=o(1)$.

It remains to bound the divisor tail beyond the forced block. For $k<h\le L$ and $d\le\sqrt{3X}$, write $g=(d,Q)$. If $d\mid n+h$ is soluble then $g\mid h$: indeed no fresh prime divides $n+h$ and $n\equiv0\pmod D$. The hyperbola bound $\tau(n+h)\le2\sum_{d\mid n+h,\,d\le\sqrt{3X}}1$ and \eqref{crt:crt} give
\begin{align*}
 \E_n\tau(n+h)
 &\le 2\sum_{\substack{d\le\sqrt{3X}\\(d,Q)\mid h}}\frac{(d,Q)}{d}
   +O(Q/\sqrt X)\\
 &\ll \tau(h)(1+\log X)+Q/\sqrt X .
\end{align*}
For the last inequality, group $d=g j$ by the divisor $g\mid h$, and bound each inner sum by $\sum_{j\le\sqrt{3X}}1/j$. The mean weighted sum for $k<h\le L$ is therefore $O((1+\log X)(k+1)2^{-k}+Q/\sqrt X)=o(1)$. The tail past $L$ has the same crude estimate used in \eqref{crt:bshort}. Finally, $\sum_{h\le k}2^{-h}\tau(n+h)$ is an integer on our class. The mean of $\dist(T_{\N}(n),\Z)+|T_B(n)-S_B|$ tends to zero, so a choice of one $n$ per $k$ proves the lemma.
\end{proof}

\begin{proof}[Proof of Theorem~\ref{crt:main}]
Suppose $uE+vS_B\in\Z$ for integers $u,v$. By the scaled-tail formula, $uT_{\N}(n)+vT_B(n)\in\Z$ for every $n$. Lemma~\ref{crt:return} forces $vS_B\in\Z$ by taking distances to the closed set $\Z$. Hence $uE\in\Z$, which gives $u=0$ by Erd\H{o}s's theorem. If $B$ is infinite, the finite-reciprocal-sum theorem gives $v=0$ as well. Clearing denominators establishes the independence of $1,E,S_B$, equivalently that of $1,S_B,E-S_B$. The irrationality of $E-S_B$ also follows directly by taking $u=-v\ne0$, including the case of finite $B$.
\end{proof}

\subsection{Examples and boundary}
Any subset of the squares satisfies the hypothesis for $s>1/2$. More generally, every powerful integer can be written as $a^2b^3$ with $b$ squarefree, by separating odd and even prime exponents. Therefore their Dirichlet series at $s>1/2$ is bounded by $\zeta(2s)\zeta(3s)$, and arbitrary deletions from this family are permitted. For infinite $B$ these examples give the joint conclusion on both sides of the partition, including arbitrary aperiodic selections.

For the support consisting of all products of two distinct primes, neither it nor its complement satisfies the sparsity assumption. Its divisor transform is $\binom{\omega(m)}{2}$. The first parity condition in the divisor block cannot be imposed throughout any progression $m\equiv a\pmod{M}$. Indeed put $g=(a,M)$ and $H=M/g$. Choose $j$ distinct primes coprime to $gM$, with product $v$, and a further prime $r$ in the reduced residue class $r\equiv (a/g)v^{-1}\pmod H$. By Dirichlet's theorem there are infinitely many choices avoiding all previous factors. Then $m=gvr$ and $\omega(m)=\omega(g)+j+1$. Taking $j=0,1,2,3$ reaches every residue modulo four, hence both parities of $\binom{\omega(m)}{2}$. This rules out the first parity step of this particular progression argument for the complete semiprime support.


\section{Periodic supports and arbitrary sparse perturbations}
\label{sec:sparse-transfer}
For an integer-valued function $c$ of bounded size write
\[
 L(c)=\sum_{m\ge1}\frac{c(m)}{2^m-1},\qquad
 g_c(t)=\sum_{b\mid t}c(b),\qquad
 U_c(n)=\sum_{h\ge1}2^{-h}g_c(n+h).
\]
Then $U_c(n)-2^nL(c)$ is an integer. Call a set $B$ sparse in this
section if $\sum_{b\in B}b^{-s}<\infty$ for some $0<s<1$.
Finite unions of such sets are sparse, by taking the largest exponent
from the finitely many summability conditions. Write $V$ for the
$\mathbb Q$-linear span of the numbers $S_B$ with sparse $B$.

\begin{lemma}\label{lem:sparsetransfer}
Let $a$ be a bounded integer-valued function. Suppose there is a set
of primes $\mathcal R$ with $\#(\mathcal R\cap[P,2P])\gg P/\log P$
for all sufficiently large $P$, for which
\begin{equation}\label{transfer-toggle}
 a(rd)=a(d)\qquad(r\in\mathcal R,\ r\nmid d).
\end{equation}
If $L(a)$ is irrational, then $L(a)\notin\mathbb Q+V$.
\end{lemma}
\begin{proof}
Clearing denominators, a contrary relation has the form
$uL(a)+L(c)\in\mathbb Z$, where $u$ is a nonzero integer and $c$ is a
bounded integer-valued function supported on one sparse set $B$.
Use the randomized Chinese remainder construction for sparse deletions
above, with all the available fresh primes taken from
$\mathcal R\cap[P,2P]$. Its proof uses only the displayed lower bound
on the number of available primes. In particular the sampling estimate
for $W_F(b)$ and the resulting nonsmooth-divisor bound are unchanged.
For $h$ up to the truncation cutoff, the smooth divisors of $n+h$
and $h$ coincide. Triangle inequality bounds the error from the remaining
divisors by $\|c\|_\infty$ times the same two counts as in the proof for
$B$. The elementary end-tail estimate also scales by $\|c\|_\infty$.
Thus the mean of $|U_c(n)-L(c)|$ on the selected progression tends to zero.

At initial shift $i$ we have $i$ distinct assigned primes having exact
valuation one in $n+i$. Every divisor of $n+i$ has a unique representation
$d\prod_{r\in I}r$ where $d$ divides the part coprime to all the assigned
primes and $I$ is a subset of them. Equation \eqref{transfer-toggle}
makes $a$ constant across the $2^i$ choices of $I$ for each $d$.
Therefore $2^i\mid g_a(n+i)$. For later shifts, $|g_a|\le
\|a\|_\infty\tau$, so the preceding divisor-tail estimates give a
vanishing mean distance from $U_a(n)$ to $\mathbb Z$. Choose a common
$n$ at every scale by summing these two mean errors.
The integers $uU_a(n)+U_c(n)$ force $L(c)\in\mathbb Z$ in the limit,
whence $uL(a)\in\mathbb Z$, a contradiction.
\end{proof}

\begin{corollary}\label{cor:togglesupp}
Let $C$ be a nonempty set whose indicator satisfies \eqref{transfer-toggle}
for a prime set with the supply in Lemma~\ref{lem:sparsetransfer}.
Whenever $A\triangle C$ is sparse, $S_A$ is irrational.
\end{corollary}
\begin{proof}
Use the same progression, taking $a=\mathbf1_C$ and $B$ empty. The
first $k$ terms of $U_a(n)$ have an integer sum, and the remaining
sum is strictly positive: for any fixed $c\in C$ there is $h>k$ with
$c\mid n+h$. Its mean tends to zero by the divisor-tail bound. Thus
one can choose $n$ for each $k$ with this positive remainder tending
to zero. Rationality $qS_C\in\mathbb Z$ would make each remainder a
multiple of $1/q$. This proves the required irrationality of $S_C$.
The signed sum $S_A-S_C$ belongs to $V$, so the lemma finishes the proof.
\end{proof}

This applies, for example, to all integers whose prime factors belong
to a set $\mathcal R$ with the stated supply, and to arbitrary sparse
changes to any such support. For $\mathcal R$ the primes congruent
to $1\bmod 4$, the unperturbed support has natural density zero:
avoiding each finite subset of the primes $3\bmod4$ gives the upper
density bound $\prod(1-1/p)$, which tends to zero by Mertens' theorem
for that progression. Its reciprocal sum diverges because it contains
the primes $1\bmod4$. Sparse changes preserve both the divergence
and density zero. Each permitted prime $p$ still divides infinitely
many composite exponents, since removing all but finitely many
$p r$ with $r\equiv1\pmod4$ would violate sparsity.

\begin{theorem}\label{thm:robust-periodic}
If $a$ is a nonzero purely periodic rational sequence, then $L(a)$
does not belong to $\mathbb Q+V$. In particular, let $C$ be an infinite
eventually periodic support and let $A\triangle C$ be sparse. Then
$S_A$ is irrational. This permits arbitrary additions and deletions
within the sparse set.
\end{theorem}
\begin{proof}
After a common denominator is cleared, $a$ is integer-valued. For its
period $M$, primes $r\equiv1\pmod M$ satisfy \eqref{transfer-toggle}
and the required count by the prime number theorem for the fixed
progression. Luca and Tachiya's theorem on periodic Lambert series \cite[Theorem~1.1]{luca}
makes $L(a)$ irrational; apply Lemma~\ref{lem:sparsetransfer}.
For an eventually periodic indicator use its periodic extension from
the tail. Its difference from $\mathbf1_A$ is supported on the sparse
symmetric difference together with finitely many elements. Its values
are $0,1,-1$, so its Lambert sum belongs to $V$.
\end{proof}

For example one may begin with all integers $1\bmod 3$ and reverse
membership on an arbitrary subset of the squares. Even without any
regularity in the selected squares, the perturbed sum is irrational.

\begin{corollary}\label{cor:rfreesparse}
Fix $e\in\{2,3,\ldots\}\cup\{\infty\}$, and let $C_e$ consist of the
positive integers free of $e$th powers of primes, with $C_\infty=\mathbb N$.
The numbers $S_{jC_e}$, $j=1,2,\ldots$, are linearly independent modulo
$\mathbb Q+V$. The same conclusion holds when $C_e$ is replaced by the
set of positive integers representable as a sum of two integer squares.
\end{corollary}
\begin{proof}
Take finitely supported integers $b_j$, not all zero, and define
$a(m)=\sum_{j\mid m}b_j\mathbf1_C(m/j)$, where $C$ is the given set.
Then $L(a)=\sum_j b_j S_{jC}$ is irrational by Duverney and Tachiya \cite[Corollary~1.2]{duverney} with
 $\ell=1$, taking $E$ to be all primes in the
$e$-free case and taking their Example~1.2 in the two-squares case.
Verify \eqref{transfer-toggle} using primes $r$ not dividing any $j$
with $b_j\ne0$. Provided $r\nmid d$, multiplying by $r$ preserves
membership of $d/j$ in $C_e$, since it introduces an exponent of one.
For the two-squares case restrict further to $r\equiv1\pmod4$: the
criterion that primes $3\bmod4$ have even exponent is then unaffected.
Each prime supply satisfies the density requirement. The lemma gives
the asserted independence of the cosets.
\end{proof}

In contrast, for the indicator of products of two distinct primes,
\eqref{transfer-toggle} fails for every $r$: take $d$ to be another
prime. The new transfer therefore does not settle the unrestricted
semiprime sum.


\section{Conditional independence of the first sixty-nine squarefree layers}
\label{sec:pratt69}
Write $\omega(m)$ for the number of distinct prime factors, with
$\omega(1)=0$. For $j\ge1$ the set $A_j$ of all products of $j$
distinct primes has sum
\[
 H_j=S_{A_j}=\sum_{k\ge1}2^{-k}\binom{\omega(k)}{j}.
\]
This follows because $\binom{\omega(k)}{j}$ counts its members dividing $k$.

\begin{theorem}\label{th:pratt69}
Assume Pratt's quantitative prime tuples conjecture
\cite[Conjecture~1.2]{pratt}. Then $1,H_1,\ldots,H_{69}$
are linearly independent over $\mathbb Q$.
\end{theorem}

\begin{proof}
For every polynomial $P$ of degree at most $69$ that is integer-valued
on the nonnegative integers put
\[
 H(P)=\sum_{k\ge1}2^{-k}P(\omega(k)),\qquad
 (\Delta P)(t)=P(t+1)-P(t).
\]
We first prove $H(P)\in\mathbb Q\Longrightarrow H(\Delta P)\in\mathbb Q$.

For $x\to\infty$ set $u=\log\log x$ and
\[
 K=\lfloor 99.8\log u\rfloor,\quad L=\lfloor 110u\rfloor,
 \quad Q=\operatorname{lcm}(1,\ldots,K),\quad J=u(\log u)^2.
\]
The prime number theorem gives $\log Q=(1+o(1))K$, so
$Q=u^{99.8+o(1)}\le u^{100}$ for sufficiently large $x$.
The $K$ distinct linear forms $1+(Q/k)n$, $1\le k\le K$, meet both
coefficient bounds and the dimension bound $K\le100\log u$ in
Pratt's conjecture. They are admissible, since the residue $n=0$ is
always available. For $p\le2K$ the local factor of their singular
series is at least $1/p\ge1/(2K)$ by admissibility. For $p>2K$ the
$K$ slopes are distinct and nonzero modulo $p$, as $p\nmid Q$ and
$1,\ldots,K$ are distinct modulo $p$. The logarithm of the local factor
$(1-K/p)(1-1/p)^{-K}$ is bounded below by $-O(K^2/p^2)$.
Multiplication thus bounds the singular series below by
$\exp(-O(K\log K))$. Pratt's formula gives at least
$x\exp(-O(u\log u))$ integers $n\le x$ for which all forms are prime.

We can impose $n>K^2$ and $\omega(nQ+k)\le J$ for every $K<k\le L$.
Indeed the number violating any latter condition is at most
\[
 L2^{-J}\sum_{m\le 2xQ}2^{\omega(m)}
 \ll L2^{-J}xQ\log(2xQ),
\]
by $2^{\omega(m)}\le\tau(m)$ and the elementary divisor sum estimate.
The logarithm of the ratio of this upper bound to $x$ is at most $-(\log2)u(\log u)^2+O(u)$.
After discarding these exceptions and the first $K^2$ integers, good
choices of $n$ remain for all sufficiently large $x$. Put $N=nQ$. For $1\le k\le K$,
the prime factor $1+nQ/k$ exceeds $k$, and is therefore coprime to $k$.
The identity $N+k=k(1+nQ/k)$ gives $\omega(N+k)=\omega(k)+1$.

For $K<k\le L$ the absolute contribution of the shifted tail is
$O_P(2^{-K}(1+J)^{69})=o(1)$ since
$99.8\log2>69$. For $k>L$ the crude bound
$\omega(N+k)\ll\log(2xQ)+\log(1+k)$ gives
\[
 \sum_{k>L}2^{-k}|P(\omega(N+k))|
 \ll_P 2^{-L}\big((\log(2xQ))^{69}+(1+L)^{69}\big)=o(1),
\]
since $\log x=\exp(u)$ and $110\log2>69$. As $K\to\infty$ this proves
\[
 \sum_{k\ge1}2^{-k}P(\omega(N+k))
 \longrightarrow\sum_{k\ge1}2^{-k}P(\omega(k)+1)
 =H(P)+H(\Delta P).
\]
If $qH(P)\in\mathbb Z$ for a positive integer $q$, then $q$ times
each shifted tail is an integer after subtracting the first $N$ terms
from $q2^NH(P)$. Its limit is an integer, giving $qH(\Delta P)\in\mathbb Z$.

A nontrivial rational relation gives a rational $H(P)$ with
$P(t)=\sum_{j=1}^d c_j\binom{t}{j}$, $1\le d\le69$ and $c_d\ne0$
after clearing denominators. By Pascal's identity, applying the implication
$d-1$ times makes $c_dH_1$ plus an integer rational.
Tao and Ter\"av\"ainen's unconditional irrationality of $H_1$
\cite[Theorem~1.3]{tt} contradicts this.
\end{proof}


\section{One integer denominator for all squarefree dilations}
For $\mathcal S=\{pr:p<r\text{ primes}\}$ write $T=S_{\mathcal S}$.
The number of elements of $\mathcal S$ dividing $m$ is
$\binom{\omega(m)}2$, so $T=\sum_{m\ge1}2^{-m}\binom{\omega(m)}2$.
\begin{theorem}\label{thm:semiprime-dilations}
There are $C>0$, integers $n_j\to\infty$ and a function $\rho:\N\to\N$ with
$1\le\rho(h)\le Ch$ such that, for each fixed $h$, $\omega(n_j+h)$ eventually equals
$\omega(h)+\rho(h)$. For every squarefree $d\ge1$, $d\mid n_j$ eventually and
\[
 \omega(n_j/d+h)\longrightarrow\omega(h)+\rho(dh)\qquad(h\text{ fixed}).
\]
If $qT\in\Z$ with $q\ge1$, then every such sequence constructed below satisfies
\begin{equation}\label{eq:semiprime-dilations}
 q\sum_{h\ge1}2^{-h}\left(\rho(dh)\omega(h)+\binom{\rho(dh)}2\right)\in\Z
 \qquad\text{for all squarefree }d\ge1 .
\end{equation}
\end{theorem}
\begin{proof}
Write $\Omega(m)$ for the number of prime factors with multiplicity. In the construction proving Theorem~1.1 of Tao and Ter\"av\"ainen \cite[Section~2, in particular axiom (i)]{tt}, for arbitrarily large $x$ one can choose $n\in[x,2x]$ divisible by $p^2$ for every prime $p\le w(x)$, where $w(x)\to\infty$, with $\Omega(n+k)\le Ck$ for all positive $k$.
The divisibility holds on the entire support of their sieve distribution, and their simultaneous bound holds on a set of positive probability in it.

For fixed $h$, eventually every prime dividing $h$ divides $n+h$. If no further prime divided $n+h$, the multiplicity bound would imply $n+h\le\max(2,h)^{Ch}$, impossible for large $n$. Hence $1\le\omega(n+h)-\omega(h)\le Ch$ eventually for that $h$. Successive subsequences and a diagonal choice yield $n_j$ and $\rho$. For fixed squarefree $d$, eventually $d^2\mid n_j$. With $m=n_j/d+h$ we have
\[
 \omega(dm)-\omega(m)=\#\{p\mid d:p\nmid m\}
 =\#\{p\mid d:p\nmid h\}=\omega(dh)-\omega(h).
\]
Applying stabilization at $dh$ proves the second assertion. Since
$\omega(m)\le\Omega(n_j+dh)\le Cdh$, the summands
$2^{-h}\binom{\omega(n_j/d+h)}2$ have a summable majorant for each fixed $d$. If $qT\in\Z$, multiplying the binary expansion by $q2^v$ and removing its finite initial portion gives
$q\sum_{h\ge1}2^{-h}\binom{\omega(v+h)}2\in\Z$ at every integer $v\ge0$.
Take $v=n_j/d$, pass to the limit by dominated convergence, and subtract $qT$. The identity $\binom{a+b}2-\binom a2=ab+\binom b2$ gives \eqref{eq:semiprime-dilations}.
\end{proof}

For example, if $\rho(dh)$ is eventually equal to a constant $c$ for some squarefree $d$, the series in \eqref{eq:semiprime-dilations} before multiplication by $q$ is
$c\sum_{h\ge1}2^{-h}\omega(h)+r$, with $r\in\Q$ and $c\ge1$. This contradicts the unconditional irrationality theorem \cite[Theorem~1.3]{tt}. More generally let $\rho(dh)$ be eventually periodic in $h$ for some squarefree $d$. Its periodic extension is nonzero, since $\rho\ge1$. Theorem~\ref{mixcube:main} shows that the corresponding periodic multiple of $\omega(h)$ has an irrational sum against $2^{-h}$. The binomial correction is eventually periodic, and the finite changes to the first term are rational. This again contradicts \eqref{eq:semiprime-dilations}. The sieve construction provides no assertion of eventual constancy or periodicity.

\section{Irrationality from agreement on quadratic intervals}

\begin{theorem}\label{thm:quadratic-agreement}
Let $A\subseteq\N$ and fix a positive integer $d$. If, for infinitely many positive integers $H$, the indicators of $A$ and $d\N$ agree at every integer in
$(H,\lceil 3H^2/4\rceil]$, then $S_A$ is irrational.
\end{theorem}
\begin{proof}
Let $w_m=(2^m-1)^{-1}$, $B=d\N$, and $J=\lceil 3H^2/4\rceil$.
The number $S_B=h_{2^d}(1)$ is irrational with irrationality exponent
$\mu(S_B)\le 2.46497868\ldots$ by Zudilin \cite[Theorem~1]{zu}. Assume $qS_A\in\Z$ for a fixed $q\ge1$. Put
\[
 R_H=S_A-\sum_{m\le H}(\ind_{m\in A}-\ind_{m\in B})w_m,
 \qquad D_H=\prod_{k\le H}\Phi_k(2),
\]
where $\Phi_k$ is the $k$th cyclotomic polynomial. The denominator of $R_H$ divides $qD_H$, since $2^m-1=\prod_{k\mid m}\Phi_k(2)$. Zudilin's Lemma~1 gives
$\log_2 D_H=(3/\pi^2+o(1))H^2$, where the logarithm here is to base two.
Agreement on the interval yields
\[
 0<|R_H-S_B|\le\sum_{m>J}w_m\le2^{1-J}.
\]
The strict inequality follows from $R_H\in\Q$ and $S_B\notin\Q$. The reduced denominators of these approximants are unbounded, since $R_H\to S_B$. For every $\nu<\pi^2/4$, the last error is eventually less than $(qD_H)^{-\nu}$ and hence less than the reduced denominator to the power $-\nu$. This forces
$\mu(S_B)\ge\pi^2/4=2.4674011\ldots$, a contradiction.
\end{proof}

Theorem~\ref{thm:quadratic-agreement} allows arbitrary behavior outside the specified intervals. For instance set $H_j=2^{3^j}$ and
$A=\bigcup_{j\ge1}((H_j,\lceil3H_j^2/4\rceil]\cap\N)$. Then $S_A$ is irrational. Both $A$ and its complement have upper natural density one, as seen respectively at the right endpoints and just before the next interval. Its exponent reciprocal sum diverges, and every prime divides infinitely many composite elements of $A$.

Taking the empty set as the comparison support gives a separate statement. If $A$ is infinite and $S_A$ is rational, then for each $\lambda>3/\pi^2$ the interval $(H,\lceil\lambda H^2\rceil]$ contains an element of $A$ for every sufficiently large $H$. Indeed the positive error after truncation is at least $1/(qD_H)$ under rationality, whereas a gap makes it at most $2^{1-\lceil\lambda H^2\rceil}$.

\section{Irrationality under sparse changes to the multiples of an integer}
\label{sec:discrepancy-denominators}
Write $S_A=\sum_{n\in A}(2^n-1)^{-1}$ and fix $d\ge1$.
The reference value for $D=d\mathbb N$ is
$S_D=h_{2^d}(1)$, where
$h_b(1)=\sum_{n\ge1}(b^n-1)^{-1}$.
Zudilin's Theorem~1 \cite{zu} gives
$\mu(S_D)\le\mu_\star=2.46497868\ldots$ for every fixed $d$.
Here $\mu$ denotes the irrationality exponent. We use this bound with the
actual discrepancy indices between the two supports. Large intervals of
quadratic length are unnecessary when their Mersenne denominators have
large common factors.

\begin{theorem}\label{thm:support-denominators}
Let $m_1<m_2<\cdots$ enumerate the symmetric difference
$A\mathbin{\triangle}D$, assumed infinite, and put
\[
 D_j^*=\operatorname{lcm}_{1\le i\le j}(2^{m_i}-1).
\]
If $\limsup_{j\to\infty} m_{j+1}/\log_2 D_j^*>\mu_\star$, then
$S_A$ is irrational. If $A\mathbin{\triangle}D$ is finite,
$S_A$ is also irrational.
\end{theorem}

\begin{proof}
A finite discrepancy changes $S_D$ by a rational number. Otherwise put
$\sigma_i=\mathbf1_A(m_i)-\mathbf1_D(m_i)\in\{-1,1\}$ and write
$R_j=\sum_{i\le j}\sigma_i/(2^{m_i}-1)$.
Its denominator divides $D_j^*$, and
\[
 |S_A-S_D-R_j|
 \le \sum_{n\ge m_{j+1}}\frac1{2^n-1}
 \le 2^{2-m_{j+1}}.
\]
If $S_A=a/q$ were rational, $a/q-R_j$ would approximate the
irrational $S_D$ with reduced denominator $b_j\le qD_j^*$.
All errors are nonzero because $S_D$ is irrational.
The denominators tend to infinity because the approximants converge
to $S_D$. Choose $\mu_\star<\nu<\nu'$ both smaller than the stated
limsup. Infinitely often $m_{j+1}>\nu'\log_2 D_j^*$. Along those
indices, $2^{2-m_{j+1}}<(qD_j^*)^{-\nu}\le b_j^{-\nu}$
eventually, contradicting the irrationality-exponent bound.
\end{proof}

\begin{corollary}\label{cor:divisor-chains}
Suppose $A\mathbin{\triangle}d\mathbb N$ is contained in a divisibility
chain $n_1\mid n_2\mid\cdots$ of distinct positive integers such that
$n_{k+1}\ge3n_k$ for all sufficiently large $k$. Then $S_A$ is irrational.
In particular one may change membership at any subset of the exponents
$1,m,m^2,\ldots$ for an arbitrary integer $m\ge3$.
\end{corollary}

\begin{proof}
For an infinite selection from the chain we have $m_i\mid m_j$ whenever
$i\le j$, so $D_j^*=2^{m_j}-1$ exactly. The ratio in
Theorem~\ref{thm:support-denominators} has lower limit at least three.
A finite selection is covered by the other part of the theorem.
\end{proof}

For a second application take changes on the dyadic chain
$\mathcal D=\{2^j:j\ge0\}$. If both the set of changes and its
complement in $\mathcal D$ are infinite, successive changed exponents
have ratio at least four infinitely often, and
Theorem~\ref{thm:support-denominators} applies. It leaves the
fully changed chain and its finite modifications undecided by this
estimate alone. Theorem~\ref{thm:robust-periodic} settles this case
as well. In fact for every $B\subseteq\mathcal D$ one has
$\sum_{n\in B}n^{-s}\le(1-2^{-s})^{-1}$ for every $s>0$, and the
multiples of $d$ form a nonempty periodic support.

In fact the sparse perturbation theorem gives joint independence:
\begin{corollary}\label{cor:q-harmonic-dyadic}
The numbers $1$, $G(1/2)=\sum_{j\ge0}(2^{2^j}-1)^{-1}$, and
$h_{2^d}(1)$ for $d=1,2,\ldots$ are linearly independent over
$\mathbb Q$.
\end{corollary}
\begin{proof}
A putative relation with finitely many nonzero rational $b_d$ gives
the periodic coefficient $a(m)=\sum_{d\mid m}b_d$. Its Lambert
sum is $\sum_d b_dh_{2^d}(1)$. If some $b_d$ is nonzero then
$a$ is nonzero by evaluation at the least such $d$. Theorem
\ref{thm:robust-periodic} excludes membership of this periodic
Lambert sum in the sum of $\mathbb Q$ and the rational span of
Lambert sums on sparse supports. Since $\mathcal D$ is sparse,
$G(1/2)=S_{\mathcal D}$ lies in that span. Hence all $b_d=0$.
The irrationality of $G(1/2)$ eliminates the remaining relation:
indeed Coons \cite{coons} proves its irrationality exponent equals two.
\end{proof}

The same argument permits many discrepancies per interval:
if all discrepancy exponents up to some cutoff divide an integer $M$
and the next discrepancy exceeds $(\mu_\star+\eta)M$
for a fixed $\eta>0$, the criterion applies provided this happens
after arbitrarily many discrepancies. Indeed $D_j^*\mid 2^M-1$.

For each fixed $d$ the class in Corollary~\ref{cor:divisor-chains} has continuum many distinct values. For $d=1$ an explicit example is
$A=\mathbb N\setminus\{3^j:j\ge0\}$ with harmonic sum
$\sum_{a\in A}a^{-1}=\infty$ and infinitely many composite exponents
divisible by each prime. Distinctness follows from strict tail
domination. For $w_n=(2^n-1)^{-1}$,
\[
 w_n-\sum_{k>n}w_k
 =\sum_{r\ge2}2^{-rn}\left(1-\frac1{2^r-1}\right)>0.
\]
At the first index where two supports differ, the difference at
that index therefore exceeds every possible cancellation from the tail.

\section{The denominators of semiprime truncations}

\begin{proposition}
Write $T_X=\sum_{a\in\mathcal S,\,a\le X}w_a$ in lowest terms with denominator $b_X$. Then $\log_2 b_X\gg X^2/\log X$, and
\[
 \lim_{X\to\infty}\frac{-\log(T-T_X)}{\log b_X}=0 .
\]
\end{proposition}
\begin{proof}
For each prime $r\in(X/4,X/2]$ with $r>3$, set $M_r=(2^r+1)/3$. Since $r\equiv1$ or $5\pmod6$, checking $2^r\pmod9$ shows $3\nmid M_r$. For each prime $\ell\mid M_r$, the multiplicative order of $2$ modulo $\ell$ divides $2r$ but does not divide $r$. Order two would force $\ell=3$, so the order is exactly $2r$. Hence $\ell$ can divide $2^a-1$ for $a\in\mathcal S$ only if $a=2r$. In the finite sum the $\ell$-power from that one denominator therefore survives reduction, and the $M_r$ for distinct $r$ are pairwise coprime. Their product divides $b_X$. The prime number theorem gives $\gg X/\log X$ such $r$, each with $\log_2 M_r\gg X$. Finally Bertrand's postulate supplies a prime $r$ with $X<2r\le 2X+O(1)$ for large $X$. Its term gives $T-T_X\ge2^{-2X-O(1)}$, proving the limit.
\end{proof}

This rules out obtaining a contradiction from the ordinary truncation error being smaller than $1/(qb_X)$ for a fixed hypothetical denominator $q$. Theorem~\ref{thm:quadratic-agreement} does not settle $T$ either. Indeed for fixed $d>1$ choose a prime $t\nmid d$. For large $X$ there is a prime $r\in(X/t,2X/t+O(1)]$ distinct from $t$ and coprime to $d$. Then $tr\in\mathcal S$ but $d\nmid tr$. For $d=1$ a prime in the same scale witnesses a disagreement. The actual semiprime support therefore differs from every fixed multiples set already on a linear scale.



\section{What remains for the two original targets}\label{sec:scope}
\paragraph{The single semiprime sum.}
For distinct finite sets of forbidden primes, the difference $T_F-T_G$
is irrational. Among the $2^{10}$ subsets of any ten specified primes
this certifies that at least $1023$ of the $1024$ sums are irrational.
It does not identify a possible exception. Theorem~\ref{th:pratt69}
includes $T_\varnothing=H_2$ under its stated conjecture.
Theorem~\ref{thm:semiprime-dilations} supplies a necessary condition for
rationality. An unconditional contradiction to that condition, or a
different unconditional proof for $T_\varnothing$ itself, is still missing.

\paragraph{An arbitrary infinite set.}
The proved sufficient hypotheses include finite harmonic mass,
bounded incidence of composite exponents, sparse changes to periodic
or prime-stable supports, and agreement on long intervals.
Corollary~\ref{cor:arbitrarysplit} applies to \emph{every} infinite
set but makes a claim about sums of its parts. Its total may be the sum
of two transcendental numbers. Strict tail domination also makes the
map $A\mapsto S_A$ injective, so only countably many supports can
correspond to algebraic numbers. Neither assertion eliminates a
particular rational candidate with unrestricted infinite support.
No unconditional proof of the full statement of problem 257 is given.

\paragraph{Provenance and scope.}
These are consolidated research notes. Standard results and external
analytic inputs are cited where used. The closest prior work includes
the prime-factor series of Tao and Ter\"av\"ainen, the periodic Lambert
series of Luca and Tachiya, and the independence results of Duverney
and Tachiya. The extensions and proofs recorded here have not undergone
independent peer review. The ten explicit transcendental sums in
Section~\ref{sec:lacunary} are constructed examples. The finite
certificate for Euler's constant and the two open targets have
different quantifiers, which are retained in their respective statements.


\begin{thebibliography}{99}
\bibitem{arthan} R. D. Arthan, \emph{The Eudoxus real numbers}.
\href{https://arxiv.org/abs/math/0405454}{arXiv:math/0405454}.
\bibitem{cdt} F. Calegari, V. Dimitrov and Y. Tang,
\emph{The linear independence of $1$, $\zeta(2)$, and $L(2,\chi_{-3})$}.
\href{https://arxiv.org/abs/2408.15403v2}{arXiv:2408.15403v2}.
\bibitem{cook} W. Cook, \emph{Irrationality criteria for Lambert subseries},
author's project and manuscript (2026).
\href{https://wcook04.github.io/plectis/maths/problems/erdos_257.html}{Problem and proof records}.
\bibitem{coons} M. Coons, \emph{On the rational approximation of the sum of the reciprocals of the Fermat numbers}.
\href{https://arxiv.org/abs/1112.5072}{arXiv:1112.5072}.
\bibitem{duverney} D. Duverney and Y. Tachiya, Refinement of the Chowla--Erd\H{o}s method and linear independence of certain Lambert series,
\emph{Forum Math.} 31 (2019), 1557--1566. \href{https://doi.org/10.1515/forum-2018-0299}{Publisher record}.
\bibitem{erdos} P. Erd\H{o}s, On arithmetical properties of Lambert series,
\emph{J. Indian Math. Soc. (N.S.)} 12 (1948), 63--66.
\bibitem{eg} P. Erd\H{o}s and R. L. Graham, \emph{Old and New Problems and Results in Combinatorial Number Theory}, Monographies de L'Enseignement Math\'ematique 28 (1980), p.~62.
\href{https://www.erdosproblems.com/257}{Problem 257 and historical references}.
\bibitem{fraenkel} A. S. Fraenkel, Iterated floor function, algebraic numbers, discrete chaos, Beatty subsequences, semigroups,
\emph{Trans. Amer. Math. Soc.} 341 (1994).
\href{https://doi.org/10.1090/S0002-9947-1994-1138949-9}{Publisher record}.
\bibitem{kimberling} C. Kimberling, Partitioning the positive integers with higher order recurrences,
\emph{Internat. J. Math. Math. Sci.} 14 (1991), 457--462.
\href{https://emis.univie.ac.at/journals/HOA/IJMMS/Volume14_3/721296.pdf}{Article}.
\bibitem{luca} F. Luca and Y. Tachiya, Irrationality of Lambert series associated with a periodic sequence,
\emph{Int. J. Number Theory} 10 (2014), 623--636. \href{https://doi.org/10.1142/S1793042113501121}{Publisher record}.
\bibitem{pratt} K. Pratt, \emph{The irrationality of a prime factor series under a prime tuples conjecture}.
\href{https://arxiv.org/abs/2409.15185}{arXiv:2409.15185}, especially Conjecture~1.2.
\bibitem{shallit} J. Shallit, \emph{An “experimental mathematics” approach to Stolarsky interspersions via automata theory}.
\href{https://arxiv.org/abs/2502.03312}{arXiv:2502.03312}.
\bibitem{tt} T. Tao and J. Ter\"av\"ainen, \emph{Quantitative correlations and some problems on prime factors of consecutive integers}.
\href{https://arxiv.org/abs/2512.01739v2}{arXiv:2512.01739v2}, especially Sections 2, 3 and 5.
\bibitem{williams} K. S. Williams, Mertens' theorem for arithmetic progressions,
\emph{J. Number Theory} 6 (1974), 353--359. \href{https://doi.org/10.1016/0022-314X(74)90032-8}{Publisher record}.
\bibitem{zu} W. Zudilin, Heine's basic transform and a permutation group for $q$-harmonic series,
\emph{Acta Arith.} 111 (2004), 153--164. \href{https://www.impan.pl/shop/publication/transaction/download/product/82435}{Article}.
\end{thebibliography}
\end{document}

